\section{Folding model and verification}
\label{app:folding-model}

\subsection{Flat-folded states}
\label{sec:problem-states}

A state represents the positions of paper regions in the plane, their orientations, and their
layer ordering. Original-sheet coordinates identify each region and preserve its connectivity to
other regions as folds change its current position. Overlapping regions can occupy the same
planar area while remaining distinct layers. Their original-sheet identities therefore matter
both when validating a move and when comparing the completed construction with the target.
Layer ordering is also a central component of computational treatments of flat-folded states
\citep{akitaya-flatfolder}.

\subsection{Supported actions}
\label{sec:problem-actions}

The simple-fold literature distinguishes operations by the layers they move
\citep{arkin-map}. Our all-layers action selects the paper on one side of a fold line throughout
the stack. Our some-layers action restricts that selection to a contiguous run from the top or
bottom. An action specifies the line in current folded coordinates, the moving side, the layer
selection, and whether the moving paper folds over or under the stationary paper. The selected
paper is reflected across the line, with the corresponding orientation and stack-order updates.
In the implemented some-layers model, a top run folds over and a bottom run folds under.


\subsection{Connectivity and tearing}
\label{app:tearing}

A move must preserve the connections between regions of the original sheet. If moving paper is
attached to stationary paper along a segment away from the proposed hinge, the operation would
separate connected paper and is rejected with a \texttt{would-tear} response. The check uses
original-sheet connectivity together with the proposed motion; overlap in the current plane alone
cannot establish whether two regions are attached. Rejection leaves the current state unchanged.

The engine also checks the action parameters and compatibility with the target crease pattern. It represents
zero-thickness flat states and does not test continuous collisions during the folding motion.

\subsection{Depth and search difficulty}

Reference depth is the number of generation steps. Search effort also depends on the
branching structure and alternative solutions. For random proposal, the probability of
selecting a legal move depends on both the legal set and the proposal distribution.

\subsection{Computational complexity}
\label{app:folding-complexity}

The sequence-generation method of \citet{akitaya-cp2seq} and the hardness proofs concern
different contributions. For a precise hardness result, \citet[Theorem 3]{akitaya-hard}
prove that deciding simple foldability of square paper with assigned creases at multiples
of $45^\circ$ is strongly NP-complete in their some-layers and all-layers models.
The input is an arbitrary crease pattern; the question is whether any permitted sequence
folds it. Strong NP-hardness persists when numerical magnitudes are polynomially bounded.

The reduction starts from \textsc{3-Partition}: divide integers into triples with a common
sum. In the orthogonal-paper construction (Theorem 1), pleating a staircase encodes choices
of integers. A bar must pass through successive cage slots; the geometry forces each group
of choices to contain three integers with the required sum. A partition therefore yields
a folding sequence, and a complete folding yields a partition. Theorem 3 embeds this
construction in square paper using additional creases that enforce the preparatory folds.
The construction has polynomial size, and a proposed folding order can be checked in
polynomial time.

For CP2Seq, each target is generated by a recorded legal sequence, so reachability is
promised rather than decided. The model must recover a sequence that also matches a
supplied terminal state, under the implemented action rules and numerical tolerances.
The theorem above does not by itself prove hardness of this promised reconstruction
problem or of our sample distribution. Its relevance is that general sequence selection
can encode combinatorial choices even when individual actions can be checked efficiently.
In an episode, the practical search burden depends on the available actions, solution
depth, and the number of unsuccessful branches explored; fold count alone does not
determine that burden.

\subsection{Comparison implementation}
\label{app:comparison-details}

The engine compares creases in original-sheet coordinates after grouping and merging collinear
intervals by assignment. Interval endpoints are matched within $2\times10^{-6}$.

The shared terminal comparator first checks layer counts. It considers direct alignment and
turnover, with turnover reversing layer order and face parity. Correspondences between vertices
of the first layer propose global alignments; a candidate alignment is accepted only if every
layer polygon and parity matches its target rank. Polygon matching uses a positional tolerance
of $2\times10^{-6}$, short-edge cleanup at $10^{-7}$, and a collinearity threshold of $10^{-9}$.

Table~\ref{tab:tolerances} collects every numerical tolerance the reported results depend on.
Each result record carries the values it was scored under, so a re-scoring can be distinguished
from a comparison at the same settings.

\begin{table}[h]
\centering
\caption{Numerical tolerances used by the reported checks.}
\label{tab:tolerances}
\begin{tabular}{llp{0.40\textwidth}}
\hline
Value & Where used & What it admits \\
\hline
$2\times10^{-6}$ & Crease comparator; terminal comparator & Interval endpoint and polygon vertex
positions, after an arbitrary global isometry. \\
$10^{-7}$ & Terminal comparator & Short-edge cleanup before polygon matching. \\
$10^{-9}$ & Terminal comparator & Collinearity threshold for vertex removal. \\
$2\times10^{-6}$ & CP edit distance & Interval endpoint agreement. \\
$2.44\times10^{-4}$ & CP edit distance & Line quantization when grouping collinear intervals. \\
$10^{-5}$ & CP edit distance & Gap below which two intervals on one line are merged. \\
\hline
\end{tabular}
\end{table}

The shared implementation represents current-plane polygons, parity, and rank. It does not
explicitly compare original-sheet region identities. Any identity-aware offline comparator must
be distinguished from this path, and the comparator version used for each reported result must
be recorded. In particular, equal shapes at different original-sheet locations can require a
stronger correspondence check than current-plane geometry alone.

The repository contains comparison and verification tests; executed test results and the exact
comparator revision must accompany any reported validation.
% EDITING NOTE: Capture fresh test output and resolve identity-aware versus shared comparator
% differences before claiming all positive and corrupted-state controls pass.

\subsection{Action rejection conditions}
\label{app:action-errors}
\begin{table}[h]
\centering
\caption{Geometric and action constraints checked by the implementation. Rejected edits preserve
 the previous state.}
\begin{tabular}{lp{.65\linewidth}}
\hline
Code & Condition \\
\hline
\texttt{would-tear} & A moving region and a stationary region share an original-sheet edge
whose current position is not on the fold line. Reflecting one side would separate the join. \\
\texttt{no-crease} & A nonempty moving selection produces no crease segment in the paper.
This condition is checked separately from tearing. \\
\texttt{nothing-to-move} & The selected half-plane contains no selected paper, or the requested
partial selection contains no layers. \\
\texttt{direction-impossible} & The action requests folding a top run under the stack or a bottom
run over it, contrary to the supported layer-selection rule. \\
\hline
\end{tabular}
\end{table}

The independent Flat-Folder wrapper uses layer-order constraints, including taco-taco,
taco-tortilla, tortilla-tortilla, and transitivity constraints
\citep{akitaya-flatfolder,flatfolder}.
It tests whether the pattern admits a flat-folded state.
